\chapter{Experimental Setup and Test Functions}
\section{Hardware and software}
\Cref{tab:system} lists the computer used and the size of the DE search. A short speed test at the start of the project was used to choose a search size (tier \ComputeTier) that fits the available time.

\begin{table}[htbp]\centering\small
\caption{Computer used and DE search size. The starting population counts toward the total number of trials.}\label{tab:system}\label{tab:budget}
\begin{minipage}[t]{.61\linewidth}\vspace{0pt}\centering
\begin{adjustbox}{max width=\linewidth}\begin{tabular}{ll}
\toprule
Item & Measured / configured value \\
\midrule
CPU & Apple M4 Pro \\
Memory (GiB) & 48 \\
Device & mps \\
Python & 3.11.14 \\
PyTorch & 2.15.0 \\
Operating system & macOS-26.7.1-arm64-arm-64bit \\
Compute tier & B \\
\bottomrule
\end{tabular}
\end{adjustbox}
\end{minipage}\hfill
\begin{minipage}[t]{.35\linewidth}\vspace{0pt}\centering
\begin{adjustbox}{max width=\linewidth}\begin{tabular}{lr}
\toprule
Setting & Value \\
\midrule
Proxy resolution & 128 \\
Training fraction & 20\% \\
Proxy epochs & 3 \\
DE population & 8 \\
DE generations & 8 \\
Budget per search & 72 \\
Scale F & 0.6 \\
Crossover CR & 0.9 \\
Search seed & 42 \\
\bottomrule
\end{tabular}
\end{adjustbox}
\end{minipage}
\end{table}

\section{Neural-network protocol}

\ArtifactTable{tab_experiment_protocol}{Training and evaluation settings for the neural-network experiments.}{tab:protocol}

\Cref{tab:protocol} lists the training settings. All choices are made on validation data; the test images are used only once, at the very end, after every recipe has been fixed. Each final recipe is trained three times with different random seeds to show how much results vary between runs. Speed (latency) is measured as the time to classify one image.

\section{Benchmark test functions}
Before using DE on the CNN, we checked it on five standard test functions~\cite{jamil2013functions}. Each has a known best value of zero. Sphere is a simple bowl; Rosenbrock is a long curved valley; Rastrigin, Ackley, and Griewank have many local traps. The tests use ten dimensions; the pictures show two.
\begin{align}
f_{\mathrm{Sphere}}(\mathbf x)&=\sum_{i=1}^{D}x_i^2,\\
f_{\mathrm{Rosenbrock}}(\mathbf x)&=\sum_{i=1}^{D-1}\left[100(x_{i+1}-x_i^2)^2+(1-x_i)^2\right],\\
f_{\mathrm{Rastrigin}}(\mathbf x)&=10D+\sum_{i=1}^{D}\left[x_i^2-10\cos(2\pi x_i)\right],\\
f_{\mathrm{Ackley}}(\mathbf x)&=-20e^{-0.2\sqrt{D^{-1}\sum_i x_i^2}}
-e^{D^{-1}\sum_i\cos(2\pi x_i)}+20+e,\\
f_{\mathrm{Griewank}}(\mathbf x)&=1+\frac{1}{4000}\sum_{i=1}^{D}x_i^2-\prod_{i=1}^{D}\cos\left(\frac{x_i}{\sqrt{i}}\right).
\end{align}
\begin{table}[htbp]\centering\small
\caption{Search range and location of the best point for each test function.}
\begin{tabular}{lll}\toprule
Function & Domain per coordinate & Minimizer\\\midrule
Sphere & $[-5.12,5.12]$ & $x_i=0$\\
Rosenbrock & $[-5,10]$ & $x_i=1$\\
Rastrigin & $[-5.12,5.12]$ & $x_i=0$\\
Ackley & $[-32.768,32.768]$ & $x_i=0$\\
Griewank & $[-600,600]$ & $x_i=0$\\\bottomrule
\end{tabular}
\end{table}

\ArtifactFigureShort{fig_benchmark_surfaces}{Two-dimensional benchmark landscapes}{The five test functions drawn in two dimensions.}{fig:surfaces}

\section{Fair optimiser comparison}
Every method gets the same budget of \BenchmarkBudget\ function evaluations and is run \BenchmarkRuns\ times with different seeds, giving \TotalBenchmarkRuns\ runs and \TotalBenchmarkEvaluations\ evaluations in total.

The population methods start from 40 evenly spread points; random search samples points independently. DE uses $F=0.6$ and $CR=0.9$. jDE starts from random $F$ and $CR$ values and keeps the ones that work. GA uses tournament selection, blend crossover, small Gaussian mutations, and elitism. PSO uses the standard settings $w=0.7298$ and $c_1=c_2=1.49618$.

Curves show the median run with its middle-50\% band; tables show mean, spread, median, and best. A rank-sum test compares DE with each other method (\Cref{tab:benchmedian,tab:benchstats}).

\section{Corruption test}
To see how the models cope with poorer photos, the seed-zero baseline and DE models were tested on the same test images with five kinds of damage: blur, noise, lower brightness and contrast, JPEG compression, and low light. Each was applied at three strengths, and the reported robustness score is the average accuracy over all of them.
